\section{Variabel, Operasi Dasar, dan Evaluasi Boolean}

\subsection{Variabel Boolean}

Variabel Boolean adalah objek matematis diskrit yang hanya dapat menyimpan salah satu dari tepat dua nilai: $0$ atau $1$. Tidak terdapat nilai antara, nilai desimal, maupun nilai negatif dalam domain variabel Boolean. Sifat biner ini secara langsung berkorespondensi dengan representasi sinyal pada rangkaian digital, di mana tegangan tinggi (\textit{high}) merepresentasikan nilai 1 dan tegangan rendah (\textit{low}) merepresentasikan nilai 0 \cite{uw_boolean}. Variabel Boolean biasanya dinotasikan dengan huruf kecil seperti $x, y, z$ atau huruf kapital $A, B, C$.

Perbedaan fundamental antara variabel Boolean dan variabel aljabar konvensional terletak pada kardinalitas domain. Variabel konvensional memiliki domain tak berhingga ($\mathbb{R}$), sedangkan variabel Boolean memiliki domain berhingga dengan kardinalitas dua. Keterbatasan ini justru menjadi kekuatan utama aljabar Boolean, karena memungkinkan enumerasi lengkap atas semua kemungkinan nilai melalui tabel kebenaran. Untuk fungsi dengan $n$ variabel Boolean, jumlah total kombinasi inputnya adalah $2^n$, sebuah angka yang meskipun tumbuh secara eksponensial, tetap berhingga dan dapat dievaluasi secara sistematis \cite{rosen2019}.

\subsection{Tiga Operasi Fundamental}

Terdapat tiga operasi dasar dalam aljabar Boolean yang menjadi fondasi bagi seluruh manipulasi ekspresi Boolean. Ketiga operasi ini saling melengkapi dan mencukupi untuk merepresentasikan fungsi Boolean apa pun (\textit{functionally complete}).

\begin{enumerate}[label=\textbf{\arabic*.}, itemsep=8pt]
  \item \textbf{Penjumlahan Boolean / OR ($+$):} Operasi biner yang menghasilkan nilai $1$ jika \emph{setidaknya satu} operand bernilai $1$, dan menghasilkan nilai $0$ hanya jika \emph{kedua} operand bernilai $0$. Operasi ini berkorespondensi dengan disjungsi ($\vee$) pada logika proposisi.

  \item \textbf{Perkalian Boolean / AND ($\cdot$):} Operasi biner yang menghasilkan nilai $1$ \emph{hanya jika kedua} operand bernilai $1$, dan menghasilkan nilai $0$ jika setidaknya satu operand bernilai $0$. Operasi ini berkorespondensi dengan konjungsi ($\wedge$) pada logika proposisi.

  \item \textbf{Komplemen / NOT ($'$ atau $\overline{\phantom{x}}$):} Operasi uner yang membalikkan nilai Boolean. Jika $x = 0$, maka $x' = 1$; jika $x = 1$, maka $x' = 0$. Operasi ini berkorespondensi dengan negasi ($\neg$) pada logika proposisi.
\end{enumerate}

Tabel~\ref{tab:operasi-dasar-boolean} menampilkan tabel kebenaran untuk ketiga operasi dasar tersebut secara lengkap.

\begin{table}[H]
\centering
\renewcommand{\arraystretch}{1.3}
\begin{tabular}{|c|c||c|c|c|}
\hline
$x$ & $y$ & $x + y$ (OR) & $x \cdot y$ (AND) & $x'$ (NOT) \\ \hline
0 & 0 & 0 & 0 & 1 \\ \hline
0 & 1 & 1 & 0 & 1 \\ \hline
1 & 0 & 1 & 0 & 0 \\ \hline
1 & 1 & 1 & 1 & 0 \\ \hline
\end{tabular}
\caption{Tabel Kebenaran Operasi Dasar Aljabar Boolean}
\label{tab:operasi-dasar-boolean}
\end{table}

\begin{catatan}
Pada tabel di atas, kolom $x'$ (NOT) hanya bergantung pada nilai $x$, bukan pada nilai $y$. Operator NOT bersifat uner dan beroperasi hanya pada satu variabel, berbeda dengan OR dan AND yang bersifat biner.
\end{catatan}

\subsection{Prioritas Operator}

Sebagaimana dalam aljabar konvensional, operator-operator Boolean memiliki urutan prioritas (\textit{precedence}) yang menentukan urutan evaluasi ketika sebuah ekspresi tidak menggunakan tanda kurung secara eksplisit. Urutan prioritas dari tertinggi ke terendah adalah sebagai berikut:

\begin{enumerate}
  \item \textbf{NOT ($'$)} --- prioritas tertinggi, dievaluasi pertama kali
  \item \textbf{AND ($\cdot$)} --- dievaluasi setelah komplemen
  \item \textbf{OR ($+$)} --- prioritas terendah, dievaluasi terakhir
\end{enumerate}

Dengan demikian, ekspresi $x + y \cdot z'$ dievaluasi sebagai $x + (y \cdot (z'))$, bukan $(x + y) \cdot (z')$ maupun bentuk lainnya. Pengetahuan tentang urutan prioritas ini sangat penting untuk menghindari ambiguitas dalam penulisan dan evaluasi ekspresi Boolean.

\subsection{Evaluasi Fungsi Boolean}

Fungsi Boolean adalah pemetaan dari $\{0,1\}^n$ ke $\{0,1\}$ yang dibentuk dari variabel Boolean, konstanta Boolean, dan operator-operator Boolean. proses evaluasi fungsi boolean dilakukan dengan mensubstitusikan nilai-nilai input ke dalam ekspresi dan menerapkan operasi-operasi sesuai urutan prioritas.

\begin{contoh}
\textbf{Evaluasi Fungsi Boolean Langkah demi Langkah}

Diberikan fungsi Boolean $F(x,y,z) = x \cdot (y + z')$. Evaluasi fungsi untuk input $(x, y, z) = (1, 0, 1)$:

\begin{enumerate}
    \item Substitusikan nilai: $F(1,0,1) = 1 \cdot (0 + 1')$
    \item Evaluasi komplemen (prioritas tertinggi): $1' = 0$
    \[ F = 1 \cdot (0 + 0) \]
    \item Evaluasi OR: $0 + 0 = 0$
    \[ F = 1 \cdot 0 \]
    \item Evaluasi AND: $1 \cdot 0 = 0$
    \[ F = 0 \]
\end{enumerate}
\end{contoh}

\begin{contoh}
\textbf{Tabel Kebenaran Lengkap untuk Fungsi Boolean}

Untuk fungsi $F(x,y,z) = x \cdot (y + z')$, tabel kebenaran lengkap dengan semua $2^3 = 8$ kombinasi input ditampilkan berikut ini.

\begin{center}
\renewcommand{\arraystretch}{1.2}
\begin{tabular}{|c|c|c||c|c|c|}
\hline
$x$ & $y$ & $z$ & $z'$ & $y + z'$ & $F = x \cdot (y + z')$ \\ \hline
0 & 0 & 0 & 1 & 1 & 0 \\ \hline
0 & 0 & 1 & 0 & 0 & 0 \\ \hline
0 & 1 & 0 & 1 & 1 & 0 \\ \hline
0 & 1 & 1 & 0 & 1 & 0 \\ \hline
1 & 0 & 0 & 1 & 1 & 1 \\ \hline
1 & 0 & 1 & 0 & 0 & 0 \\ \hline
1 & 1 & 0 & 1 & 1 & 1 \\ \hline
1 & 1 & 1 & 0 & 1 & 1 \\ \hline
\end{tabular}
\captionsetup{hypcap=false}
\captionof{table}{Tabel Kebenaran untuk $F(x,y,z) = x \cdot (y + z')$}
\label{tab:contoh-evaluasi}
\end{center}
\end{contoh}

Proses evaluasi pada tabel kebenaran merupakan metode yang bersifat \textit{brute force} --- setiap kombinasi input dievaluasi satu per satu. Meskipun metode ini selalu menghasilkan jawaban yang benar, jumlah baris yang harus dievaluasi tumbuh secara eksponensial terhadap jumlah variabel. Untuk fungsi dengan 4 variabel diperlukan 16 baris, untuk 5 variabel diperlukan 32 baris, dan seterusnya. Keterbatasan skalabilitas inilah yang memotivasi pengembangan teknik-teknik penyederhanaan aljabar dan metode grafis seperti Peta Karnaugh yang akan dibahas pada bab berikutnya \cite{rosen2019}.
